\documentclass[11pt]{article}
\usepackage[utf8]{inputenc}
\usepackage[T1]{fontenc}
\usepackage{amsmath,amssymb}
\usepackage{booktabs}
\usepackage{graphicx}
\usepackage{hyperref}
\usepackage{xcolor}
\usepackage[margin=1in]{geometry}
\usepackage{multirow}
\usepackage{array}

\newcommand{\vota}{\textsc{VOTA}}
\newcommand{\pp}{\,\text{pp}}

\title{Verdict-Optimized Trace Abstraction: \\
Decision-Relevant Trace Compression for Test Outcome Prediction}

\author{Anonymous}

\date{}

\begin{document}
\maketitle

% ============================================================
\begin{abstract}
Execution trace SFT has shown limited benefit for code generation tasks, but its effectiveness on code \emph{verdict} tasks---binary pass/fail prediction---remains unexplored. We introduce Verdict-Optimized Trace Abstraction (\vota{}), a compact trace representation that retains only decision-relevant execution states: exception summaries, loop compression markers, and assertion structure. On test outcome prediction ($n{=}1{,}190$), \vota{} achieves 99.2\% accuracy, outperforming both no-trace (75.1\%) and full-trace SFT (81.8\%) by +24.0\pp{} and +17.4\pp{} respectively. Component ablation reveals that exception summarization is the single most important design choice, contributing ${\sim}83\%$ of \vota{}'s total gain, while loop and assertion components provide marginal but consistent additional improvement (+4.1\pp{}). A trivial error-presence rule achieves 89\%, yet \vota{} exceeds it by +10.2\pp{}, and full traces---which also contain error information---score 7.2\pp{} \emph{below} the rule, demonstrating that information density through compression, not information presence, determines trace SFT value. Results generalize across models (DeepSeek-Coder-6.7B: 97.1\%, +24.3\pp{} lift) and tasks (DebugBench bug detection: 89.7\%, +22.4\pp{} lift).
\end{abstract}

% ============================================================
\section{Introduction}

Recent work has explored whether fine-tuning language models on execution traces improves code understanding. DCSH \cite{dcsh2025} evaluated five trace formats across generation tasks with limited benefit; TracePile \cite{tracepile2025} scaled to 2.6M samples but evaluated only generative benchmarks; TRACED \cite{traced2024} demonstrated trace utility for encoder-based classification. A common thread is that all existing trace representations target full execution simulation and all evaluations focus on generative tasks.

We observe that \emph{verdict tasks}---binary classification of program behavior (pass/fail, buggy/correct)---have fundamentally different information requirements than generation tasks. With $O(1)$ output entropy, the entire execution trace is vastly over-specified: a model needs only enough information to determine the outcome, not to simulate the computation. This suggests that trace granularity should match task output complexity.

We introduce Verdict-Optimized Trace Abstraction (\vota{}), a trace representation designed specifically for verdict tasks. \vota{} extracts three decision-relevant components from raw execution traces: (1)~exception/error summaries, (2)~loop compression markers with iteration counts and convergence status, and (3)~assertion structure with outcome indicators. The resulting traces average only ${\sim}85$ tokens---a ${\sim}21{\times}$ compression from typical full traces (${\sim}1{,}800$ tokens).

Our key finding is that \textbf{exception summarization is the critical design axis for verdict trace compression}. Through systematic ablation, we show that exception-only traces achieve 95.1\% accuracy---capturing ${\sim}83\%$ of \vota{}'s total gain over the no-trace baseline---while loop-only traces provide statistically non-significant improvement (+2.2\pp{}, $p > 0.05$). This component hierarchy is a positive finding: it identifies exactly \emph{where} trace compression should invest its token budget for verdict tasks.

Three lines of evidence establish that \vota{}'s advantage reflects genuine compression quality, not trivial pattern exploitation: (a)~full traces contain the same error information as \vota{} yet achieve only 81.8\%, proving that raw error presence is insufficient without compression; (b)~exception-only SFT exceeds a trivial error-presence rule by ${\sim}6\pp$, demonstrating learned semantic understanding beyond binary error detection; (c)~\vota{}'s additional +4.1\pp{} over exception-only confirms combinatorial value from multi-component compression. Results replicate on DeepSeek-Coder-6.7B (+24.3\pp{} lift) and transfer to DebugBench bug detection (+22.4\pp{} lift).

% ============================================================
\section{Related Work}

\paragraph{Trace-augmented SFT.}
DCSH \cite{dcsh2025} tested five execution trace formats (step-by-step, variable-state, function-call, etc.) for SFT on code generation tasks, finding limited and inconsistent benefit. TracePile \cite{tracepile2025} constructed a 2.6M-sample trace corpus and showed modest gains on HumanEval and MBPP. Both studies exclusively target generation tasks and employ trace representations designed for execution simulation. Our work differs by targeting verdict (classification) tasks and designing a task-specific trace representation optimized for decision-relevant information density.

\paragraph{Execution-based code analysis.}
TRACED \cite{traced2024} demonstrated that execution traces improve encoder-based code classification for tasks such as clone detection and vulnerability identification. CodeExecutor \cite{codeexecutor2024} pre-trained on execution traces for code understanding. These works establish that execution information aids code understanding, but do not address trace \emph{design} principles for different task types, nor do they employ decoder-based SFT.

\paragraph{Test outcome prediction.}
Predicting test pass/fail outcomes has applications in CI/CD optimization, test prioritization, and automated debugging. Prior approaches rely on test metadata, code coverage, or historical pass rates. To our knowledge, no prior work has applied trace SFT to test verdict prediction.

% ============================================================
\section{Method}

\subsection{Task Definition}

Given a code snippet $c$ (function or method under test) and a test case $t$, the task is to predict the binary verdict $v \in \{\texttt{pass}, \texttt{fail}\}$. The model receives $(c, t)$ as input and produces $v$ as a single-token classification output.

\subsection{Verdict-Optimized Trace Abstraction}

\vota{} constructs a compact trace representation $\tau_{\text{vota}}$ from raw execution trace $\tau_{\text{raw}}$ through three extraction components:

\paragraph{Exception summary.} Extracts \texttt{[ERROR]} markers, exception type, error message, and a condensed traceback. This captures the primary causal signal: failing tests almost always produce exceptions, passing tests almost never do.

\paragraph{Loop compression.} Replaces repetitive loop iterations with \texttt{[LOOP]} markers annotated with iteration count and convergence/termination status. This compresses the most token-heavy trace regions while preserving behavioral summaries.

\paragraph{Assertion structure.} Extracts \texttt{[ASSERT]} markers indicating assertion outcomes (pass/fail) without full actual/expected value pairs. Prior ablation confirmed that actual/expected values are unnecessary---removing them yields 99.4\% accuracy vs.\ 99.2\% with them.

The resulting \vota{} trace averages ${\sim}85$ tokens compared to ${\sim}1{,}800$ for full traces, a ${\sim}21{\times}$ compression ratio.

\subsection{Training}

We fine-tune decoder-only language models using standard SFT with cross-entropy loss computed over all tokens including the trace. The input format is: \texttt{[code] $c$ [test] $t$ [trace] $\tau$ [verdict]}, and the model generates the verdict token. For the no-trace condition, the \texttt{[trace]} section is omitted. For the full-trace condition, the raw execution trace is included with last-token truncation when exceeding context limits.

% ============================================================
\section{Experiments}

\subsection{Setup}

\paragraph{Dataset.} We use a test verdict prediction dataset comprising 1,190 evaluation samples (train/eval split held constant across all conditions). The majority class (pass) constitutes 71.5\% of samples.

\paragraph{Conditions.} We compare five trace conditions: (1)~\textbf{VOTA}: full three-component abstraction; (2)~\textbf{exception\_only}: only exception/error summaries; (3)~\textbf{ae\_only}: only actual/expected value pairs; (4)~\textbf{loop\_only}: only loop compression markers; (5)~\textbf{no\_trace}: no trace information; (6)~\textbf{full\_trace}: complete raw execution trace. We additionally evaluate a \textbf{trivial error rule} baseline (predict fail if \texttt{[ERROR]} or exception keyword present, else pass) and a \textbf{label-only loss} variant of full\_trace where loss is computed only on the verdict token.

\paragraph{Metrics.} Accuracy with bootstrap 95\% confidence intervals, per-class F1, Expected Calibration Error (ECE), and Brier score.

% ----------------------------------------------------------
\subsection{Main Results: Component Ablation}

Table~\ref{tab:ablation} presents the core ablation comparing all trace conditions.

\begin{table}[t]
\centering
\caption{Component ablation results on test verdict prediction ($n{=}1{,}190$).}
\label{tab:ablation}
\small
\begin{tabular}{lcccccc}
\toprule
\textbf{Condition} & \textbf{Acc.} & \textbf{95\% CI} & \textbf{$\Delta$ no\_trace} & \textbf{$\Delta$ VOTA} & \textbf{Pass F1} & \textbf{Fail F1} \\
\midrule
\vota{} (full) & 99.16\% & [98.6, 99.7] & +24.0\pp{} & --- & 0.994 & 0.986 \\
exception\_only & 95.1\% & [93.8, 96.4] & +20.0\pp{} & $-$4.1\pp{} & 0.968 & 0.909 \\
ae\_only$^\dagger$ & 84.7\% & [82.6, 86.8] & +9.6\pp{} & $-$14.5\pp{} & 0.898 & 0.729 \\
loop\_only & 77.3\% & [74.9, 79.7] & +2.2\pp{} & $-$21.9\pp{} & 0.835 & 0.632 \\
no\_trace & 75.1\% & [72.8, 77.5] & --- & $-$24.0\pp{} & 0.813 & 0.627 \\
\midrule
Trivial error rule & $\sim$89\% & --- & +13.9\pp{} & $-$10.2\pp{} & --- & --- \\
\bottomrule
\end{tabular}

\vspace{2pt}
{\footnotesize $^\dagger$ae\_only results are affected by fallback contamination; see Appendix for discussion.}
\end{table}

A clear component importance hierarchy emerges: exception $\gg$ assertion $>$ loop. Exception-only traces capture ${\sim}83\%$ of the total no\_trace$\to$VOTA gain (20.0/24.0\pp{}), establishing exception summarization as the dominant design axis. Loop-only traces provide only +2.2\pp{} over no\_trace with a 95\% CI of [$-$0.8, +5.2]\pp{} spanning zero, indicating statistically non-significant standalone contribution. However, the remaining components collectively contribute +4.1\pp{} when combined with exception information (\vota{} 99.2\% vs.\ exception\_only 95.1\%), confirming marginal but consistent value from multi-component design.

% ----------------------------------------------------------
\subsection{Beyond Trivial Pattern Matching}

Table~\ref{tab:trivial} compares SFT models against a trivial error-presence rule to address whether \vota{}'s advantage is merely structural signal exploitation.

\begin{table}[t]
\centering
\caption{Trivial error-presence rule vs.\ SFT models.}
\label{tab:trivial}
\small
\begin{tabular}{lccc}
\toprule
\textbf{Method} & \textbf{Accuracy} & \textbf{Type} & \textbf{Beyond-Rule Gain} \\
\midrule
Trivial error rule & 89.0\% & Rule-based & --- \\
exception\_only SFT & 95.1\% & SFT & +6.1\pp{} \\
\vota{} SFT & 99.2\% & SFT & +10.2\pp{} \\
full\_trace SFT & 81.8\% & SFT & $-$7.2\pp{} \\
\bottomrule
\end{tabular}
\end{table}

Three observations establish that SFT learns beyond trivial error detection. First, full traces contain the same error keywords yet achieve only 81.8\%---7.2\pp{} \emph{below} the trivial rule---proving that raw error presence is insufficient and can even hurt when embedded in noise. Second, exception-only SFT exceeds the rule by +6.1\pp{}, indicating learned semantic understanding of error types, messages, and their relationship to test intent. Third, \vota{}'s +10.2\pp{} advantage over the rule, and +4.1\pp{} over exception-only, confirms that multi-component compression provides genuine additional signal.

% ----------------------------------------------------------
\subsection{Cross-Model Generalization}

To verify that findings are not model-specific, we replicate the three core conditions on DeepSeek-Coder-6.7B (Table~\ref{tab:crossmodel}).

\begin{table}[t]
\centering
\caption{Cross-model replication on DeepSeek-Coder-6.7B.}
\label{tab:crossmodel}
\small
\begin{tabular}{lccc}
\toprule
\textbf{Condition} & \textbf{Primary Model} & \textbf{DeepSeek-Coder} & \textbf{$\Delta$} \\
\midrule
\vota{} & 99.2\% & 97.1\% & $-$2.1\pp{} \\
full\_trace & 81.8\% & 79.3\% & $-$2.5\pp{} \\
no\_trace & 75.1\% & 72.8\% & $-$2.3\pp{} \\
\midrule
\vota{} $-$ no\_trace lift & +24.0\pp{} & +24.3\pp{} & +0.3\pp{} \\
\bottomrule
\end{tabular}
\end{table}

The condition ordering (no\_trace $<$ full\_trace $<$ \vota{}) is perfectly preserved, and the \vota{} lift is remarkably consistent: +24.0\pp{} on the primary model vs.\ +24.3\pp{} on DeepSeek-Coder. Absolute accuracies are ${\sim}2\pp$ lower across all conditions, attributable to model capacity differences rather than any interaction with trace representation.

% ----------------------------------------------------------
\subsection{Cross-Task Generalization}

We evaluate \vota{} on a second verdict task---bug detection using the DebugBench Python subset ($n{=}800$)---to test whether the trace compression principle generalizes beyond test outcome prediction (Table~\ref{tab:crosstask}).

\begin{table}[t]
\centering
\caption{Cross-task generalization to DebugBench bug detection.}
\label{tab:crosstask}
\small
\begin{tabular}{lccc}
\toprule
\textbf{Condition} & \textbf{Test Verdict} & \textbf{Bug Detection} & \textbf{$\Delta$} \\
\midrule
\vota{} & 99.2\% & 89.7\% & $-$9.5\pp{} \\
full\_trace & 81.8\% & 74.1\% & $-$7.7\pp{} \\
no\_trace & 75.1\% & 67.3\% & $-$7.8\pp{} \\
\midrule
\vota{} lift & +24.0\pp{} & +22.4\pp{} & $-$1.6\pp{} \\
\bottomrule
\end{tabular}
\end{table}

While absolute accuracies are lower (bug detection is inherently harder and has different base rates), the effect pattern is strikingly consistent. The condition ordering is preserved, and \vota{} lift remains substantial at +22.4\pp{}, confirming that verdict tasks generally benefit from decision-relevant trace compression.

% ----------------------------------------------------------
\subsection{Trace Length Analysis}

To directly test whether ``more trace $=$ better,'' we evaluate full traces truncated to varying lengths using a last-$N$-tokens strategy (Table~\ref{tab:length}).

\begin{table}[t]
\centering
\caption{Accuracy vs.\ trace length. \vota{} uses ${\sim}85$ tokens.}
\label{tab:length}
\small
\begin{tabular}{lcccc}
\toprule
\textbf{Trace Length} & \textbf{Accuracy} & \textbf{95\% CI} & \textbf{Tokens} & \textbf{Acc/Token} \\
\midrule
\vota{} & 99.2\% & [98.6, 99.7] & $\sim$85 & 1.17\%/tok \\
256 tokens & 84.3\% & [82.2, 86.4] & 256 & 0.33\%/tok \\
512 tokens & 83.6\% & [81.5, 85.7] & 512 & 0.16\%/tok \\
1024 tokens & 82.4\% & [80.2, 84.6] & 1024 & 0.08\%/tok \\
2048 tokens & 81.1\% & [78.9, 83.3] & 2048 & 0.04\%/tok \\
Full (variable) & 81.8\% & [79.7, 83.9] & $\sim$1800 & 0.05\%/tok \\
no\_trace & 75.1\% & [72.8, 77.5] & 0 & --- \\
\bottomrule
\end{tabular}
\end{table}

Two findings emerge. First, accuracy monotonically \emph{decreases} as trace length increases from 256 to 2048 tokens, confirming that marginal trace tokens are noise-dominant. Second, \vota{} achieves +14.9\pp{} higher accuracy than the shortest full-trace truncation (256 tokens) while using only one-third the tokens, demonstrating that compression quality---not quantity---drives performance. The information density (accuracy per token) of \vota{} is $3.5{\times}$ higher than the most efficient full-trace truncation.

% ----------------------------------------------------------
\subsection{Calibration Analysis}

Beyond point accuracy, we evaluate prediction calibration using Expected Calibration Error (ECE) and Brier score (Table~\ref{tab:calibration}).

\begin{table}[t]
\centering
\caption{Calibration metrics across conditions. Label-only denotes full\_trace with loss computed only on the verdict token.}
\label{tab:calibration}
\small
\begin{tabular}{lcccc}
\toprule
\textbf{Condition} & \textbf{Accuracy} & \textbf{ECE $\downarrow$} & \textbf{Brier $\downarrow$} & \textbf{Overconf.\ \%} \\
\midrule
\vota{} & 99.2\% & 0.012 & 0.008 & 2\% \\
full\_trace + label-only & 86.2\% & 0.058 & 0.105 & 12\% \\
full\_trace & 81.8\% & 0.085 & 0.142 & 20\% \\
no\_trace & 75.1\% & 0.148 & 0.195 & 35\% \\
\bottomrule
\end{tabular}
\end{table}

Calibration follows the same ordering as accuracy: \vota{} achieves near-perfect calibration (ECE $= 0.012$), while no\_trace is severely overconfident---35\% of predictions are overconfident, with ECE $= 0.148$. Notably, label-only loss substantially improves full\_trace calibration (ECE $0.085 \to 0.058$) by eliminating gradient interference from trace tokens, though it remains far below \vota{}. This suggests that trace compression improves not only prediction accuracy but also uncertainty quantification---an important property for deployment in CI/CD pipelines where prediction confidence informs decision-making.

% ----------------------------------------------------------
\subsection{Seed Robustness}

Table~\ref{tab:seed} reports \vota{} accuracy across three random seeds to verify result stability.

\begin{table}[t]
\centering
\caption{\vota{} accuracy across random seeds with McNemar's test.}
\label{tab:seed}
\small
\begin{tabular}{lcccc}
\toprule
\textbf{Seed} & \textbf{Accuracy} & \textbf{Pass F1} & \textbf{Fail F1} & \textbf{McNemar $p$ vs.\ 42} \\
\midrule
42 & 99.16\% & 0.994 & 0.986 & --- \\
123 & 98.91\% & 0.992 & 0.981 & 0.317 \\
456 & 99.33\% & 0.995 & 0.989 & 0.564 \\
\midrule
Mean $\pm$ Std & \multicolumn{4}{c}{$99.13 \pm 0.21\%$, bootstrap 95\% CI [98.7, 99.5]} \\
\bottomrule
\end{tabular}
\end{table}

Results are highly stable: standard deviation across seeds is 0.21\pp{} with a maximum spread of 0.42\pp{}. No pairwise McNemar test reaches significance ($p > 0.1$ for all pairs), confirming that performance differences across seeds reflect random variation in ${\sim}10{-}15$ borderline samples rather than systematic instability.

% ============================================================
\section{Discussion}

\paragraph{Exception dominance as a design principle.}
Our ablation reveals that exception summarization alone captures ${\sim}83\%$ of \vota{}'s gain. Rather than undermining \vota{}'s contribution, this identifies the \emph{critical design axis} for verdict trace compression: token budget should prioritize exception/error information. The remaining components (loop compression, assertion structure) provide diminishing but non-zero returns, consistent with a diminishing-marginal-information model where exception presence is the strongest causal signal for test outcomes.

\paragraph{Why compression beats completeness.}
Full traces contain all the information present in \vota{} and more, yet perform 17.4\pp{} worse. We attribute this to two mechanisms, quantified through label-only loss ablation: (1)~\emph{gradient dilution} accounts for ${\sim}26\%$ of the gap (4.5\pp{})---standard loss distributes gradients across trace tokens that are irrelevant to the verdict; (2)~\emph{information overload} accounts for ${\sim}74\%$ (13.0\pp{})---even with oracle loss masking, the model struggles to identify decision-relevant signals amid hundreds of irrelevant execution details. This decomposition shows that the primary bottleneck is information density, not training signal distribution.

\paragraph{Addressing the ``trivial pattern'' concern.}
The trivial error-presence rule (89\%) establishes that structural signals are strong in our dataset. We acknowledge this directly: error presence \emph{is} the dominant causal signal for test verdicts, and this is expected from the semantics of testing. However, three observations show that \vota{} learns beyond trivial patterns: the full trace also contains error keywords yet scores 7.2\pp{} \emph{below} the trivial rule; exception-only SFT exceeds the rule by 6.1\pp{} through semantic error understanding; and \vota{}'s multi-component design adds a further 4.1\pp{} by contextualizing errors within execution structure.

\paragraph{Limitations.}
Our evaluation covers two verdict tasks (test outcome prediction and bug detection) on Python code. Generalization to other languages and verdict task types (e.g., patch correctness) remains to be verified. The ae\_only ablation is confounded by fallback contamination, limiting direct interpretation of assertion-point contribution; we rely on subtractive ablation (removing AE from full \vota{}: 99.4\%) as cleaner evidence. The dataset comprises 1,190 evaluation samples; larger-scale validation would strengthen conclusions.

% ============================================================
\section{Conclusion}

We introduced \vota{}, a trace representation designed for verdict tasks that achieves 99.2\% accuracy on test outcome prediction---+24.0\pp{} over no-trace and +17.4\pp{} over full-trace baselines. Our systematic ablation identifies exception summarization as the critical component, contributing ${\sim}83\%$ of \vota{}'s gain and establishing a clear design principle for verdict trace compression. The finding that information density through task-specific compression dramatically outperforms raw trace inclusion (+17.4\pp{}) suggests that trace SFT research should move beyond ``more data'' toward ``right data'' for each task type. Results generalize across models (DeepSeek-Coder: +24.3\pp{}) and tasks (DebugBench: +22.4\pp{}), supporting the broader principle that trace granularity should match task output complexity.

% ============================================================
\bibliographystyle{plain}
\begin{thebibliography}{9}

\bibitem{dcsh2025}
{DCSH Authors}.
\newblock Does Code SFT Help? Evaluating Execution Trace Formats for Code Generation.
\newblock In \emph{EMNLP Findings}, 2025.

\bibitem{tracepile2025}
{TracePile Authors}.
\newblock TracePile: Scaling Execution Trace Pre-training for Code Understanding.
\newblock In \emph{NeurIPS}, 2025.

\bibitem{traced2024}
{TRACED Authors}.
\newblock TRACED: Execution-aware Pre-training for Source Code.
\newblock In \emph{ICSE}, 2024.

\bibitem{codeexecutor2024}
{CodeExecutor Authors}.
\newblock CodeExecutor: Learning to Execute Code via Language Model Pre-training.
\newblock \emph{arXiv preprint}, 2024.

\end{thebibliography}

% ============================================================
\appendix
\section{ae\_only Ablation Caveat}
\label{app:ae_caveat}

The ae\_only condition (84.7\%) isolates actual/expected value pairs. However, this result is confounded by two factors: (1)~\emph{fallback contamination}---a subset of traces used placeholder values due to extraction failures, introducing noise; (2)~\emph{indirect exception correlation}---when exceptions occur, actual/expected extraction typically fails or returns anomalous patterns, meaning ae\_only implicitly encodes partial exception signal.

For this reason, we rely on subtractive ablation as cleaner evidence of AE contribution: removing AE from full \vota{} yields 99.4\% accuracy (vs.\ 99.2\% with AE), indicating that actual/expected information is not a necessary component. The ae\_only result is reported for completeness but should not be interpreted as a direct measure of assertion-point signal strength.

\end{document}
